\documentclass[
	fontsize=12pt,
	paper=a4,
	numbers=noenddot,
	listof=totoc,
	listof=entryprefix,
	listof=nochaptergap,
	bibliography=totoc,
	parskip=half,
	openany
]{scrbook}

\usepackage{cmap}


\usepackage[utf8]{inputenc}
\usepackage[T1]{fontenc}
\usepackage{lmodern}
\usepackage{ngerman}


\usepackage{scrhack}
\usepackage[table]{xcolor}
\usepackage{chngcntr}
\usepackage{enumitem}


\usepackage{amsthm}


\usepackage{lipsum}
\usepackage{unicode-math}


\usepackage{mathtools}


\usepackage[
	colorlinks=false,
	linkcolor=black,
	citecolor=black,
	filecolor=black,
	urlcolor=black,
	bookmarks=true,
	bookmarksopen=true,
	bookmarksopenlevel=3,
	bookmarksnumbered,
	plainpages=false,
	pdfpagelabels=true,
	hyperfootnotes
]{hyperref}


\renewcommand{\familydefault}{\sfdefault}
\usepackage[onehalfspacing]{setspace}
\raggedbottom

\usepackage{anyfontsize}

\addtokomafont{chapter}{\sffamily\large\bfseries}
\addtokomafont{section}{\sffamily\normalsize\bfseries}
\addtokomafont{subsection}{\sffamily\normalsize\mdseries}
\addtokomafont{caption}{\sffamily\normalsize\mdseries}

\usepackage{relsize}
\renewcommand*{\UrlFont}{\ttfamily\smaller\relax}


\usepackage[
	bindingoffset=0cm,
	inner=2.5cm,
	outer=2.5cm,
	top=3cm,
	bottom=2cm
]{geometry}

\usepackage[
	headsepline,
	plainheadsepline
]{scrlayer-scrpage}
\clearpairofpagestyles
\addtokomafont{pagehead}{\normalfont}
\ohead*{\thepage}
\ihead*{\leftmark}


\usepackage{graphicx}
\graphicspath{{pictures/}}
\usepackage{float}

\renewcommand{\thefigure}{\arabic{figure}}
\counterwithout{figure}{chapter}

\usepackage{wrapfig}


\usepackage{booktabs}
\usepackage{multirow}
\usepackage{tabularx}

\renewcommand{\thetable}{\arabic{table}}
\counterwithout{table}{chapter}


\usepackage{listings}
\usepackage{beramono}

\renewcommand{\lstlistlistingname}{List of Code Listings}
\renewcommand{\lstlistingname}{Code Listing}
\newcommand{\listoflolentryname}{\lstlistingname}

\definecolor{codegreen}{rgb}{0,0.6,0}
\definecolor{codegray}{rgb}{0.5,0.5,0.5}
\definecolor{codepurple}{rgb}{0.5,0,0.33}
\definecolor{codepurblue}{rgb}{0.16,0.0,1.0}
\definecolor{backcolour}{rgb}{0.95,0.95,0.92}

\lstdefinestyle{codestyle}{
	backgroundcolor=\color{backcolour},
	commentstyle=\color{codegreen},
	keywordstyle=\bfseries\color{codepurple},
	numberstyle=\tiny\color{codegray},
	stringstyle=\color{codepurblue},
	basicstyle=\scriptsize\ttfamily,
	breakatwhitespace=false,
	breaklines=true,
	captionpos=b,
	keepspaces=true,
	numbers=left,
	numbersep=5pt,
	showspaces=false,
	showstringspaces=false,
	showtabs=false,
	tabsize=2,
	escapeinside={(*@}{@*)}
}
\lstset{style=codestyle, numberbychapter=false}



\KOMAoptions{toc=chapterentrydotfill}
\addtokomafont{chapterentry}{\normalfont}
\setuptoc{toc}{totoc}

\DeclareTOCStyleEntry[beforeskip=0cm]{chapter}{chapter}
\DeclareTOCStyleEntry[beforeskip=0cm]{section}{section}
\DeclareTOCStyleEntry[beforeskip=0cm]{default}{subsection}

\BeforeStartingTOC[lof]{\def\autodot{:}}
\BeforeStartingTOC[lot]{\def\autodot{:}}
\BeforeStartingTOC[lol]{\def\autodot{:}}


\iffalse
\usepackage{csquotes}
\setlength\bibitemsep{.5\baselineskip}
\setcounter{biburlnumpenalty}{9000}
\setcounter{biburllcpenalty}{9000}
\setcounter{biburlucpenalty}{9000}

\fi

\usepackage[printonlyused]{acronym}


\usepackage[title,titletoc]{appendix}

\newcommand{\appendixchapter}[1]{
\cleardoublepage
\pagenumbering{arabic}
\renewcommand{\thepage}{\thechapter-\arabic{page}}
\chapter{#1}
}

\newcommand{\monthlyreport}[2]{
\section{#1}
\centering
\includegraphics[trim=55 35 55 35,clip,width=1\textwidth]{#2}
\clearpage
}



\newtheorem{beispiel}{Beispiel}

\newtheorem{these}{These}

\newtheorem{definition}{Definition}


\usepackage[nameinlink, noabbrev]{cleveref}


\lstset{
	extendedchars=true,
	inputencoding=utf8,
	literate=
		{ä}{{\"a}}1
		{ö}{{\"o}}1
		{ü}{{\"u}}1
		{ß}{{\ss}}1
		{Ä}{{\"A}}1
		{Ö}{{\"O}}1
		{Ü}{{\"U}}1
}

\begin{document}

\begin{titlepage}
	\pagestyle{empty}

	\begin{flushright}
		\begin{figure}[ht]
			\flushright
			\includegraphics[height=2cm]{../../for_latex/hfu}
		\end{figure}
	\end{flushright}

	\begin{center}
		\vspace{3cm}

		{\fontsize{22}{22} \selectfont \textbf{Blatt 6}}\\[5mm]
		{\fontsize{18}{18} \selectfont Web-Security und Intrusion Prevention}

		\vspace{12cm}

		{\fontsize{14}{14} \selectfont Luis Staudt}
	\end{center}
\end{titlepage}


\chapter*{Aufgabe 1: Die Webanwendung}

Es laufen zwei Flask-Anwendungen: das Gästebuch aus Blatt 5 (Port 5000, mit der
dort gezeigten XSS-Lücke) und eine Online-Banking-Anwendung (Port 5001).
Benutzer der Bank sind \texttt{alice} (alice123), \texttt{bob} (bob123) und
\texttt{eve} (eve123), jeweils mit einem Kontostand. Nach dem Login setzt die
Bank ein ungeschütztes Sitzungscookie \texttt{session\_id} ohne SameSite- oder
HttpOnly-Flag.

\chapter*{Aufgabe 2: Die Banking-Seite verstehen}

Eine Überweisung von Alice an Bob über 10 EUR zeigt in der Entwicklerkonsole den
zugehörigen Aufruf. Die Überweisung läuft über einen einfachen GET-Request:

\begin{lstlisting}[language=bash]
GET /transfer?to=<empfaenger>&amount=<betrag>   (Cookie: session_id=...)
\end{lstlisting}

Benötigt werden somit nur die Parameter \texttt{to} und \texttt{amount}; die
Authentifizierung erfolgt allein über das automatisch mitgesendete
Sitzungscookie. Genau das macht die Anwendung anfällig für CSRF.

\begin{lstlisting}[language=bash]
curl -b "session_id=..." "http://127.0.0.1:5001/transfer?to=bob&amount=10"
-> Kontostand: 1000 -> 990
\end{lstlisting}

\chapter*{Aufgabe 3: CSRF über die XSS-Lücke im Gästebuch}

Überwiesen werden sollen 500 EUR an Eve. Dazu wird ein \texttt{<img>}-Tag ins
Gästebuch eingetragen, dessen \texttt{src} auf den Transfer-Endpunkt der Bank
zeigt. Das Laden eines Bildes bleibt dabei unauffällig:

\begin{lstlisting}[language=HTML]
Schoenes Bild! <img src="http://127.0.0.1:5001/transfer?to=eve&amount=500">
\end{lstlisting}

Ist Alice bei der Bank eingeloggt und ruft das Gästebuch auf, versucht ihr
Browser das vermeintliche Bild zu laden und sendet dabei ihr Banking-Cookie
automatisch mit. Die Überweisung wird im Kontext von Alice ausgeführt:

\begin{lstlisting}[language=bash]
Kontostand Alice: 1000 -> 500   (500 EUR an Eve ueberwiesen)
\end{lstlisting}

\section*{Sec-Fetch-Header der gefälschten Anfrage}

Serverseitig wurden die Header der eingehenden \texttt{/transfer}-Anfrage
protokolliert:

\begin{lstlisting}[language=bash]
Sec-Fetch-Site: same-site
Sec-Fetch-Mode: no-cors
Sec-Fetch-Dest: image
Referer:        http://127.0.0.1:5000/
Cookie:         session_id=...; secretcookie=itsasecret; ...
\end{lstlisting}

\texttt{Sec-Fetch-Site} trägt den Wert \texttt{same-site}, weil beide
Anwendungen denselben Host \texttt{127.0.0.1} nutzen; der Port gehört nicht zur
Site. Der Browser sendet das Cookie deshalb trotz des SameSite-Standards (Lax)
mit. Bei verschiedenen Sites würde Lax den Angriff bereits erschweren.
\texttt{Sec-Fetch-Dest: image} weist aus, dass die Anfrage von einem
\texttt{<img>}-Tag stammt, und \texttt{Referer} zeigt auf das Gästebuch.

\section*{Verhalten bei POST-basierter Überweisung}

Bei Verwendung von POST liesse sich der Angriff nicht mehr über ein
\texttt{<img>}-Tag auslösen, da Bilder ausschliesslich GET-Anfragen erzeugen.
Möglich bliebe er dennoch: Erforderlich wäre ein automatisch abgesendetes
HTML-Formular mit \verb|method="post"|, das per JavaScript submittet wird. Über
die XSS-Lücke lässt sich ein solches Formular ebenfalls einschleusen. POST allein
schützt somit nicht vor CSRF.

\chapter*{Aufgabe 4: Absicherung mit CSRF-Token}

Die Bank wird um CSRF-Tokens erweitert. Beim Login wird pro Sitzung ein
zufälliges Token erzeugt und serverseitig gespeichert. Das Transfer-Formular
enthält das Token als verstecktes Feld, und \texttt{/transfer} wird auf POST
umgestellt und prüft das Token:

\begin{lstlisting}[language=Python]
def set_session(user):
    sessions[session] = {"user": user, "csrf": get_random_string(32)}

@app.route("/transfer", methods=["POST"])
def transfer():
    s = get_session(request.cookies.get("session_id"))
    if not s: return redirect("/")
    if request.form.get("csrf_token") != s["csrf"]:
        abort(403)
    ...
\end{lstlisting}

Zusätzlich wird das Sitzungscookie mit \texttt{httponly=True} und
\verb|samesite="Strict"| gesetzt.

Die Tests bestätigen die Wirkung:

\begin{lstlisting}[language=bash]
# alter <img>-CSRF (GET) -> Methode nicht mehr erlaubt
GET  /transfer?to=eve&amount=500          -> 405
# gefaelschter POST ohne gueltiges Token  -> abgewiesen
POST /transfer  (ohne csrf_token)         -> 403
# legitime Ueberweisung mit Token aus dem Dashboard -> funktioniert
POST /transfer  (csrf_token=..., to=bob, amount=50) -> 302, Kontostand 1000 -> 950
\end{lstlisting}

Da ein Angreifer das sitzungsgebundene Token nicht erraten kann, wird die
gefälschte Überweisung mit 403 abgewiesen. Legitime Überweisungen aus dem
Dashboard funktionieren weiterhin.


\end{document}
